\section{Vivado/XSim and FPGA bring-up}

Vivado is a common next step after GHDL validation: it can simulate with XSim,
then synthesize/implement the same RTL on a Xilinx FPGA. XSim simulation plays
the same role here as the GHDL flow used earlier in the project (compile RTL
plus testbench, run a behavioural simulation, check the signature) -- the two
flows are interchangeable verification steps that consume the same RTL and the
same \texttt{tests/isa/<name>.sig} golden files. For pre-Vivado interactive
exploration of the pipeline in a browser (schematic view plus drag-and-drop
simulation, with real firmware loaded via \texttt{--mem-init}), see
Section~\ref{sec:digitaljs} (DigitalJS interactive simulation). DigitalJS is a
useful visual companion before committing to a Vivado project setup, but it is
not signature-checked and does not replace XSim as a verification step.

\paragraph{Automated project setup}
The full Vivado bring-up flow -- project creation, source import, VHDL-2008
file-type override, top selection, generics, and board XDC attachment -- is
captured by a single Tcl driver, \texttt{scripts/vivado\_project.tcl}, and
documented end-to-end in \texttt{docs/vivado\_instructions.md} (XSim
simulation, synthesis/implementation, and on-board ILA validation). The
sections below summarise what that script does and show the relevant Tcl
snippets so the manual steps remain reproducible if you prefer GUI clicks
over Tcl.

\begin{lstlisting}[style=minrv]
# Open project in the GUI, ready to launch sim/synth:
vivado -source scripts/vivado_project.tcl -tclargs full_core basys3
# Non-interactive (batch) creation:
vivado -mode batch -source scripts/vivado_project.tcl -tclargs full_core zedboard
\end{lstlisting}

\subsection{Vivado simulation (XSim) checklist}
The Tcl driver \texttt{scripts/vivado\_project.tcl} performs every step of
this checklist; the manual list below is mainly useful for understanding what
the script does (and what to tweak in the GUI if you ever set things up by
hand).
\begin{enumerate}[leftmargin=*]
  \item Build the program hex images (example target: \texttt{tests/isa/full\_core.c}; expected signature: \texttt{tests/isa/full\_core.sig}):
\begin{lstlisting}[style=minrv]
./scripts/build_prog.sh tests/isa/full_core.c full_core
\end{lstlisting}
  \item Create the project and import sources via the Tcl driver:
\begin{lstlisting}[style=minrv]
vivado -source scripts/vivado_project.tcl -tclargs full_core basys3
\end{lstlisting}
  Internally this adds every \texttt{rtl/*.vhd} as a design source and every
  \texttt{sim/*.vhd} as a simulation source (corresponding Tcl excerpts shown
  below in Section~\ref{sec:vivado-sources}).
  \item Force VHDL-2008 on the imported sources (Vivado defaults to VHDL-93). The driver does this automatically (see the snippet in Section~\ref{sec:vivado-2008}).
  \item Set the simulation top and point the testbench generics at the generated hex files and the expected signature file (the driver does this too; see Section~\ref{sec:vivado-tb-generics}).
  \item Run behavioral simulation (\texttt{launch\_simulation}) and look for \texttt{Note: TEST PASSED}. The driver sets the simulation run time to \texttt{all}, so the simulation runs until the testbench stops it (Vivado's default of 1000\,ns ends before \texttt{full\_core} passes at 28335\,ns).
\end{enumerate}

\subsection{Vivado source organization tip (exact lists)}
\label{sec:vivado-sources}
The documentation includes two exact source lists: one for synthesis (design
top \texttt{rtl/fpga\_top.vhd} -- a thin wrapper around \texttt{rtl/top.vhd}
that exposes only \texttt{clk} / \texttt{reset} on the pad ring), one for
simulation (simulation top \texttt{sim/tb\_minrv32.vhd}, which drives
\texttt{top} directly so the testbench still sees the observation ports).

In practice you do not type these lists into the GUI; the Tcl driver imports
them with two \texttt{glob} calls -- \texttt{rtl/*.vhd} into the design
fileset, \texttt{sim/*.vhd} into \texttt{sim\_1} -- and then selects the
correct top for each fileset:
\begin{lstlisting}[style=minrv,language=tcl]
# scripts/vivado_project.tcl (excerpt)
set rtl_files [glob -nocomplain [file join $repo rtl *.vhd]]
add_files -norecurse $rtl_files
set sim_files [glob -nocomplain [file join $repo sim *.vhd]]
add_files -fileset sim_1 -norecurse $sim_files

set_property top fpga_top   [current_fileset]
set_property top tb_minrv32 [get_filesets sim_1]
\end{lstlisting}
The exact file lists below are useful as a manual reference: anything pulled
in by the \texttt{glob} above appears here.

\paragraph{Synthesis sources}
\begin{lstlisting}[style=minrv]
rtl/alu.vhd
rtl/branch_unit.vhd
rtl/controller_pipe.vhd
rtl/core_pkg.vhd
rtl/mem_init_pkg.vhd
rtl/decode_stage.vhd
rtl/dmem.vhd
rtl/execute_stage.vhd
rtl/extend.vhd
rtl/fetch_stage.vhd
rtl/hazard_unit.vhd
rtl/imem_dp.vhd
rtl/load_extend_unit.vhd
rtl/maindec.vhd
rtl/aludec.vhd
rtl/mem_access_unit.vhd
rtl/memory_stage.vhd
rtl/mux2.vhd
rtl/mux3.vhd
rtl/pipe_reg.vhd
rtl/regfile.vhd
rtl/riscvpipe.vhd
rtl/datapath_pipe.vhd
rtl/top.vhd
rtl/fpga_top.vhd
rtl/writeback_stage.vhd
\end{lstlisting}

\paragraph{Note}
\texttt{top.vhd} instantiates \texttt{imem\_dp} (dual-port BRAM) for instruction
memory and for data-side reads from the text region (\texttt{.rodata} loads).
Port B can be left disabled if you want a logically single-port instruction
memory. \texttt{fpga\_top.vhd} is the Vivado synthesis top: it instantiates
\texttt{top} and XOR-reduces the five observation output ports into a
single 1-bit \texttt{led} so the FPGA pad ring only sees \texttt{clk} /
\texttt{reset} / \texttt{led} (3 IOs instead of 72). The XOR fold is
deliberate: it forces every observation bit to reach a pad, so Vivado
cannot prune the datapath as dead code during backwards optimisation. The
fold is XORed with the MSB of a free-running 24-bit counter (about 4.8\,Hz at
80\,MHz), so the LED blinks visibly whenever the core clock is running.

\paragraph{Core clock}
The core does not run at the board clock. \texttt{fpga\_top.vhd} contains an
MMCM (Mixed-Mode Clock Manager, the clock generator block inside 7-series
FPGAs) that makes the core clock:
core clock $=$ board clock $\times$ \texttt{MMCM\_MULT} $/$
(\texttt{MMCM\_DIVCLK} $\times$ \texttt{MMCM\_OUT\_DIV}).
The Tcl driver computes these generics from the core clock you ask for with
\texttt{--clk-mhz} (default 80\,MHz on Basys3 and ZedBoard) and picks the highest
frequency that is not above it. The board XDC constrains only the 100\,MHz
board clock; Vivado derives the core clock constraint from it, so the timing
report checks the real core clock. The core is held in reset until the MMCM
output is stable. At 80\,MHz the worst setup slack is only about 0.17\,ns on
both boards (Section~\ref{sec:vivado-results}); if a change makes it negative,
use 75\,MHz, as below.
\begin{lstlisting}[style=minrv]
vivado -mode batch -source scripts/vivado_project.tcl \
       -tclargs full_core basys3 --clk-mhz 75
\end{lstlisting}

\paragraph{Simulation sources (with optional trace wrappers)}
\begin{lstlisting}[style=minrv]
rtl/alu.vhd
rtl/branch_unit.vhd
rtl/controller_pipe.vhd
rtl/core_pkg.vhd
rtl/mem_init_pkg.vhd
rtl/decode_stage.vhd
rtl/dmem.vhd
rtl/execute_stage.vhd
rtl/extend.vhd
rtl/fetch_stage.vhd
rtl/hazard_unit.vhd
rtl/imem_dp.vhd
rtl/load_extend_unit.vhd
rtl/maindec.vhd
rtl/aludec.vhd
rtl/mem_access_unit.vhd
rtl/memory_stage.vhd
rtl/mux2.vhd
rtl/mux3.vhd
rtl/pipe_reg.vhd
rtl/regfile.vhd
rtl/riscvpipe.vhd
rtl/datapath_pipe.vhd
sim/riscvpipe_trace.vhd
rtl/top.vhd
sim/top_trace.vhd
sim/tb_minrv32.vhd
rtl/writeback_stage.vhd
\end{lstlisting}

\subsection{VHDL-2008 enable (Tcl snippet)}
\label{sec:vivado-2008}
\begin{lstlisting}[style=minrv,language=tcl]
set repo <repo_root>
foreach f [glob -nocomplain $repo/rtl/*.vhd] {
  set_property file_type {VHDL 2008} [get_files $f]
}
foreach f [glob -nocomplain $repo/sim/*.vhd] {
  set_property file_type {VHDL 2008} [get_files $f]
}
\end{lstlisting}

\subsection{Pointing the testbench at hex/signature files}
\label{sec:vivado-tb-generics}
Vivado can set testbench generics on the \texttt{sim\_1} fileset. The documented
example configures IMEM/DMEM init files plus an expected signature file:
\begin{lstlisting}[style=minrv,language=tcl]
set_property top tb_minrv32 [get_filesets sim_1]
set repo      <repo_root>
set build_dir [file join $repo build]
set_property generic \
  "IMEM_INIT_FILE=\"$build_dir/full_core.text.hex\" \
   DMEM_INIT_FILE=\"$build_dir/full_core.data.hex\" \
   EXPECTED_FILE=\"$repo/tests/isa/full_core.sig\" \
   SIGNATURE_WORDS=10 \
   MAX_CYCLES=5000 \
   HALT_ON_PASS=true" \
  [get_filesets sim_1]
\end{lstlisting}

\noindent
\texttt{SIGNATURE\_WORDS} must match the number of 32-bit words in the
\texttt{.sig} file (one per line). For example, \texttt{full\_core.sig} uses 10
words. The Tcl driver \texttt{scripts/vivado\_project.tcl} computes this
automatically by counting non-blank lines in the \texttt{.sig} file.

\subsection{Optional: per-instruction trace wrappers}
For a more didactic simulation transcript, the project includes simulation-only
wrappers that print decoded assembly and key control choices. In Vivado, add
\texttt{sim/top\_trace.vhd} and \texttt{sim/riscvpipe\_trace.vhd} to simulation
sources; for synthesis, omit them so the synthesizable RTL remains unchanged.

When the hazard unit takes action (stall/flush/forward), the trace also prints
a \texttt{HAZARD:} line with the stall/flush decisions and the forwarding
selects:
\begin{lstlisting}[style=minrv]
HAZARD: StallF=stall StallD=stall FlushD=keep FlushE=flush FwdA=rf FwdB=rf
\end{lstlisting}
\noindent
\texttt{FwdA/FwdB} select \texttt{rf}, \texttt{M}, or \texttt{W}.

\subsection{Expected warnings}
Because the simple memories mask addresses down to their depth, Vivado may warn
that higher address bits are unused. Similar ``unused'' warnings can appear for
helpers that only consume low address bits for byte strobes and alignment. The
documentation calls these warnings expected for this style of educational memory model.

\subsection{BRAM initialization: simulation vs synthesis}
The memory models share a single init mechanism: an impure
\texttt{init\_mem()} function in \texttt{rtl/mem\_init\_pkg.vhd} reads a
plain-text hex file at \emph{elaboration time} and returns a pre-populated
\texttt{word\_mem\_t} array.
\begin{itemize}[leftmargin=*]
  \item Simulation (GHDL/XSim): the function runs in the simulator's elaboration phase and pre-loads the array before time 0.
  \item Synthesis (Vivado): the same code path executes in Vivado's elaborator -- TEXTIO is supported at elaboration -- and the resulting array is baked into the inferred BRAM's \texttt{INIT\_xx} parameters in the bitstream. No tool-specific \texttt{INIT\_FILE} attribute is needed.
\end{itemize}
Provide absolute paths to the hex files via the top-level
\texttt{IMEM\_INIT\_FILE} / \texttt{DMEM\_INIT\_FILE} generics (relative paths
fail because the elaborator runs from Vivado's project working directory, not
the repo root).

\subsection{Synthesis top generics (init files and sizes)}
The synthesis top is \texttt{fpga\_top}. Its generics are
\texttt{IMEM\_WORDS}, \texttt{DMEM\_WORDS}, \texttt{IMEM\_INIT\_FILE},
\texttt{DMEM\_INIT\_FILE}, \texttt{FORWARDING\_ENABLE}, and the clock
settings \texttt{CLK\_IN\_PERIOD\_PS}, \texttt{MMCM\_MULT},
\texttt{MMCM\_DIVCLK}, and \texttt{MMCM\_OUT\_DIV} (the
\texttt{DMEM\_BASE\_ADDR} split is fixed inside \texttt{rtl/top.vhd} at
\texttt{0x2000\_0000} and is not exposed on the synthesis top). The Tcl
driver sets all of them except \texttt{FORWARDING\_ENABLE} on the design
fileset; for an 80\,MHz core clock (the Basys3 default):
\begin{lstlisting}[style=minrv,language=tcl]
set repo     <repo_root>
set build    [file join $repo build]
set_property generic \
  "IMEM_WORDS=2048 DMEM_WORDS=2048 \
   IMEM_INIT_FILE=$build/full_core.text.hex \
   DMEM_INIT_FILE=$build/full_core.data.hex \
   CLK_IN_PERIOD_PS=10000 \
   MMCM_MULT=12 MMCM_DIVCLK=1 MMCM_OUT_DIV=15" \
  [current_fileset]
\end{lstlisting}

\subsection{On-chip debugging (ILA)}
For FPGA bring-up, Vivado's Integrated Logic Analyzer (ILA) provides a
transaction-level view of execution by probing internal nets. The
\texttt{MARK\_DEBUG=true} attribute is already applied in \texttt{rtl/top.vhd}
to the nets of interest (\texttt{pc}, \texttt{instr}, \texttt{dmem\_addr\_i},
\texttt{dmem\_wdata\_i}, \texttt{dmem\_wstrb\_i}, \texttt{dmem\_write\_en\_i},
\texttt{dmem\_read\_en\_i}), so they appear directly in the ``Set Up Debug''
wizard without further RTL edits.

\subsection{Validate a program via the signature window (ILA capture)}
\label{sec:vivado-ila-signature}
The same signature idea used in simulation is used on FPGA: rather than reading
BRAM over JTAG, trigger an ILA on stores into the signature window and read the
captured store words back from the ILA buffer.
\begin{itemize}[leftmargin=*]
  \item DMEM base: \texttt{0x2000\_0000}
  \item Signature base: \texttt{0x2000\_0100} (DMEM offset \texttt{0x100})
  \item Trigger: \texttt{dmem\_write\_en\_i == 1} AND \texttt{dmem\_addr\_i == 32'h2000\_01XX}
  \item Capture: \texttt{dmem\_addr\_i} and \texttt{dmem\_wdata\_i}; depth 4096
        samples is ample (\texttt{full\_core.sig} has 10 words).
\end{itemize}

Hardware Manager Tcl to run the capture and dump the samples to CSV:
\begin{lstlisting}[style=minrv,language=tcl]
set ila [lindex [get_hw_ilas] 0]
run_hw_ila    $ila
wait_on_hw_ila $ila
write_hw_ila_data -csv_file signature_capture.csv \
                  [upload_hw_ila_data $ila]
\end{lstlisting}

Sort the captured rows by \texttt{dmem\_addr\_i} and compare the
\texttt{dmem\_wdata\_i} column against \texttt{tests/isa/<test>.sig} (one
32-bit word per line).

\subsection{Implementation results}
\label{sec:vivado-results}
Table~\ref{tab:vivado-results} lists the post-implementation results on
ZedBoard (\texttt{xc7z020clg484-1}) at the default 80\,MHz core clock, measured
with Vivado 2024.2, the default implementation strategy, and the
\texttt{full\_core} program. The ILA column is for an ILA on all 134
\texttt{MARK\_DEBUG} bits, with a depth of 4096 samples and capture control
enabled. Other RTL revisions, Vivado versions, devices, or ILA settings give
somewhat different numbers.

\paragraph{Published results}
Papers about MinRV32 report the results of the version that existed when they
were written. The FPGA results of the BCI conference paper were obtained with
an earlier version in which the core was clocked directly by the 100\,MHz board
oscillator while the timing constraint specified 80\,MHz, so its timing results
(WNS $+0.112$\,ns without and $+0.269$\,ns with the ILA) were not checked at the
clock that the core actually received. The current version fixes this with the
MMCM. Its resource usage is almost the same (1845 against 1842 LUTs without the
ILA), its estimated power includes the MMCM (0.261 against 0.162\,W without the
ILA), and its timing is checked at the real core clock. The
\texttt{docs/vivado\_instructions.md} file lists the differences in detail.

\begin{table}[htbp]
\centering
\caption{Post-implementation results on ZedBoard at 80\,MHz.}
\label{tab:vivado-results}
\begin{tabularx}{\linewidth}{@{}X r r r@{}}
\toprule
Resource / result & Available & Without ILA & With ILA \\
\midrule
Slice LUTs                  & 53200  & 1845 (3.47\%) & 3778 (7.10\%) \\
Slice registers (FF)        & 106400 & 1523 (1.43\%) & 4800 (4.51\%) \\
Occupied slices             & 13300  & 875 (6.58\%)  & 2053 (15.44\%) \\
F7 muxes                    & 26600  & 261           & 280 \\
Block RAM tiles (RAMB36)    & 140    & 4 (2.86\%)    & 19 (13.57\%) \\
DSP slices                  & 220    & 0             & 0 \\
Bonded IOB                  & 200    & 3             & 3 \\
Global clock buffers (BUFG) & 32     & 2             & 3 \\
MMCM                        & 4      & 1             & 1 \\
\midrule
WNS / WHS (ns)              & ---    & $+0.174$ / $+0.156$ & $+0.240$ / $+0.065$ \\
Total on-chip power (est.)  & ---    & 0.261\,W      & 0.289\,W \\
\quad Dynamic               & ---    & 0.152\,W      & 0.179\,W \\
\quad\quad of which MMCM    & ---    & 0.116\,W      & 0.116\,W \\
\quad Device static         & ---    & 0.108\,W      & 0.110\,W \\
\bottomrule
\end{tabularx}
\end{table}

The two BUFGs without the ILA drive the core clock and the MMCM feedback; the
ILA adds a third one for the JTAG clock of its debug hub, and it samples on the
core clock. The power is the vectorless estimate of Vivado; most of the dynamic
power is the MMCM, a fixed cost of the clock generation, and the processor
itself accounts for about 0.036\,W.

\paragraph{Critical path}
The slowest paths start at the output of the data-memory block RAM
(2.45\,ns after the clock edge) and resolve a branch that depends on the loaded
value: the IMEM/DMEM read-data mux, the load extension and the Writeback result
mux, the forwarding mux of the Execute stage, and the branch comparison produce
\texttt{PCSrcE}. \texttt{PCSrcE} then drives the next-PC adder and the PC
register, and \texttt{FlushD}/\texttt{FlushE}, which reach the synchronous reset
inputs of the Fetch/Decode and Decode/Execute registers. Without the ILA the
slowest endpoint is such a reset input (11 logic levels), and the PC register is
only 0.027\,ns behind it (19 logic levels, 9 of them \texttt{CARRY4}); with the
ILA, the PC register is the slowest endpoint. More than half of the delay is
routing, so the slack depends on placement.

\paragraph{MMCM or PLL}
Each 7-series clock tile has one MMCM and one PLL, and either can make 80\,MHz
from 100\,MHz. With a PLL at the same 1200\,MHz internal oscillator
($\times 12 / 15$), the estimated power was slightly higher (0.121\,W against
0.116\,W) and the design missed timing by 0.178\,ns, only because the core logic
was placed differently. With a PLL at 800\,MHz ($\times 8 / 10$), the total power
fell to 0.242\,W, but the clock uncertainty grew from 0.069 to 0.082\,ns and the
WNS fell to $+0.050$\,ns. The power depends on the oscillator frequency and not
on the type of block, so the MMCM is kept.
